% % \section{Related Works}
% % \label{sec:related_works}

% % % Although RAGs augment LLMs with desired characteristics like proficiency in domain-specific knowledge, personalization and data security, the interest in them has been more recent, which is reflected in its body of work and discussion in the research community. 
% % Given the growing research interest and industrial adoption for RAG and LLM based applications, in this section, we discuss prior related works. 

% % \textbf{RAG optimizations:} Prior works like FlashRAG ~\cite{jin2024flashrag} present a comprehensive framework allowing the users to choose from a diverse set of components in each stage including different indexing schemes, databases, etc. The works PipeRAG \cite{jiang2024piperag}, \cite{lu2024turborag}, \cite{zhang2024accelerating}, EdgeRAG~\cite{EdgeRAG} and \cite{ragcacheefficientknowledgecaching} present system level insights to RAG with the first two focusing on the algorithm-system co-design and the latter three proposing a caching scheme based on the documents accessed. While these solutions provide system-level insight into RAG pipeline, they are mainly targeted towards enterprise- and server-grade computing. There also exist works that present custom accelerator design and processing in-memory solutions for accelerating RAG pipeline like \cite{jiang2023chameleon} and \cite{qin2024robust}. These works show good speedup and energy efficiency numbers but do not fit our use-case considering the overheads they would bring in the integration with edge devices. These custom hardware proposals require substantial time and effort to make the solution more general for the rapidly changing generative applications. In contrast, we put forth a software-based solution which can be flexibly implemented across different platforms. Works like~\cite{huang2024toward} take a different route to solve this problem by doing RAG computations in the cloud datacenters. While this approach is more general and applicable to different devices, the concern of data security can be an issue. 

% % \textbf{LLM optimizations:} To provide the user with more accurate results, the enterprises have proposed bigger models like LLaMA 3.1~\cite{touvron2023llama}, Mixtral 8x7B~\cite{jiang2024mixtral}, GPT3~\cite{gpt3}, etc. However, as the models grow bigger, the difficulty of deploying them on edge devices grows exponentially. There are works that talk about challenges of deploying big LLMs  on edge devices~\cite{10.1145/emperical} with limited memory and compute capacity. In particular, existing body of works have proposed various optimizations for LLM execution on devices that have small memory capacities, by  selectively scheduling the subsequent layers on GPU and keeping the remaining model on the CPU storage~\cite{aminabadi2022deepspeedinferenceenablingefficient, sheng2023flexgenhighthroughputgenerativeinference}. Researchers have also explored the accuracy vs compute tradeoff space so as to find the right quantization\cite{quant, quant1} method to drop the less useful data and compute while still achieving a satisfactory accuracy. Although these optimizations have been proven useful to cut down the power budget, decrease latency and help in user's data security, these are only useful in generating response from the LLM faster and do not focus on integrating user's context in the generation output. Also, for a problem which presents as many diverse architectural aspects as RAG, the software solution needs to cater to the diverse architectures that could be running them. 

% % \textbf{Hardware resource management:} There have been works in the direction of resource-management for CPUs~\cite{gong2022graphite, \Cref, nair2024parallelization} which attempts to overlap computation and memory accesses, and better use mutltiple cores. However, these works are not directly applicable to RAG application on edge devices. 

% \section{Related Work} 
% \label{sec:related_works} %Cyan
% Recent research and industrial adoption underscore the importance of improving both RAG- and LLM-based applications. Several previous works~\cite{cacheblend, hermes, rago} investigated optimizations for server-scale RAG pipelines, including FlashRAG~\cite{jin2024flashrag}, which offers a modular framework that spans diverse indexing and database choices, PipeRAG~\cite{jiang2024piperag}, TurboRAG~\cite{lu2024turborag}, and others~\cite{zhang2024accelerating,EdgeRAG,ragcacheefficientknowledgecaching} that present system-level insights, often targeting datacenter deployments. Some proposals enhance RAG performance through specialized near storage~\cite{RAGAS, reis}, or in-memory computation~\cite{jiang2023chameleon,qin2024robust, HETERag, accelRAG, drex}, delivering notable speedups and energy efficiency, but complicating integration with resource-limited devices. For scenarios where data privacy is paramount, offloading computation to the cloud, as in \cite{huang2024toward}, can introduce confidentiality concerns. Hence, we focus on a platform-flexible software solution that suits edge settings while addressing memory and compute constraints.

% Large LLMs such as LLaMA~\cite{touvron2023llama}, Mixtral~\cite{jiang2024mixtral}, and GPT3~\cite{gpt3} have sparked research into model and resource optimization techniques for edge deployments~\cite{10.1145/emperical}. Various approaches schedule heavy layers on GPUs and offload others to CPU storage~\cite{aminabadi2022deepspeedinferenceenablingefficient, sheng2023flexgenhighthroughputgenerativeinference}, or use carefully tuned quantization schemes~\cite{quant, quant1} to reduce memory footprint with minimal accuracy loss. Such methods are critical for running large models in low-power devices, but they center on speeding up the core generative model rather than incorporating local context and retrieval. A broader software approach is needed to account for RAG’s distinct pipeline stages and the wide variety of device architectures in which they operate.

% Finally, a body of work on hardware resource management for CPUs~\cite{gong2022graphite, nair2024parallelization} explores better concurrency and compute–memory overlap, yet these solutions are not directly tailored to RAG’s unique combination of encoding, retrieval, and generation on edge platforms. Our work bridges these gaps by proposing a system-level resource-aware pipeline that addresses challenges ranging from CPU/ GPU partitioning and memory usage to overall QoS in personal computing devices.

%%%%%%%%%%%%%%%%%%%%%%%%%%%%%%%%%%%%%%%%%%%%%%%%%%%%%%%%%%%%
\section{Related Work}
\label{sec:related_works}

\paragraph{RAG Pipeline Optimization}
Prior systems such as FlashRAG~\cite{jin2024flashrag}, PipeRAG~\cite{jiang2024piperag}, TurboRAG~\cite{lu2024turborag}, and RAGCache~\cite{ragcacheefficientknowledgecaching} provide modular frameworks and system-level optimizations for RAG, but primarily target datacenter deployments with abundant GPU resources. EdgeRAG~\cite{EdgeRAG} addresses edge scenarios but relies on on-demand embedding without systematic CPU--GPU orchestration. Hardware accelerators for RAG~\cite{jiang2023chameleon, qin2024robust, reis, HETERag} achieve strong speedups but complicate integration with resource-limited devices.

\noindent\textbf{Edge ML Systems.}
Research on edge LLM deployment focuses on layer offloading~\cite{aminabadi2022deepspeedinferenceenablingefficient, sheng2023flexgenhighthroughputgenerativeinference} and quantization~\cite{quant, quant1} to reduce memory footprint. These methods accelerate the generative model but do not address the full RAG pipeline's encoding and retrieval stages.

\noindent\textbf{Heterogeneous Resource Management.}
Work on CPU resource management~\cite{gong2022graphite, nair2024parallelization} explores compute--memory overlap but is not tailored to RAG's unique stage characteristics. MaestroRAG bridges these gaps with characterization-driven, fine-grained orchestration designed specifically for edge RAG deployment.

% % \section{Related Works}
% % \label{sec:related_works}

% % % Although RAGs augment LLMs with desired characteristics like proficiency in domain-specific knowledge, personalization and data security, the interest in them has been more recent, which is reflected in its body of work and discussion in the research community. 
% % Given the growing research interest and industrial adoption for RAG and LLM based applications, in this section, we discuss prior related works. 

% % \textbf{RAG optimizations:} Prior works like FlashRAG ~\cite{jin2024flashrag} present a comprehensive framework allowing the users to choose from a diverse set of components in each stage including different indexing schemes, databases, etc. The works PipeRAG \cite{jiang2024piperag}, \cite{lu2024turborag}, \cite{zhang2024accelerating}, EdgeRAG~\cite{EdgeRAG} and \cite{ragcacheefficientknowledgecaching} present system level insights to RAG with the first two focusing on the algorithm-system co-design and the latter three proposing a caching scheme based on the documents accessed. While these solutions provide system-level insight into RAG pipeline, they are mainly targeted towards enterprise- and server-grade computing. There also exist works that present custom accelerator design and processing in-memory solutions for accelerating RAG pipeline like \cite{jiang2023chameleon} and \cite{qin2024robust}. These works show good speedup and energy efficiency numbers but do not fit our use-case considering the overheads they would bring in the integration with edge devices. These custom hardware proposals require substantial time and effort to make the solution more general for the rapidly changing generative applications. In contrast, we put forth a software-based solution which can be flexibly implemented across different platforms. Works like~\cite{huang2024toward} take a different route to solve this problem by doing RAG computations in the cloud datacenters. While this approach is more general and applicable to different devices, the concern of data security can be an issue. 

% % \textbf{LLM optimizations:} To provide the user with more accurate results, the enterprises have proposed bigger models like LLaMA 3.1~\cite{touvron2023llama}, Mixtral 8x7B~\cite{jiang2024mixtral}, GPT3~\cite{gpt3}, etc. However, as the models grow bigger, the difficulty of deploying them on edge devices grows exponentially. There are works that talk about challenges of deploying big LLMs  on edge devices~\cite{10.1145/emperical} with limited memory and compute capacity. In particular, existing body of works have proposed various optimizations for LLM execution on devices that have small memory capacities, by  selectively scheduling the subsequent layers on GPU and keeping the remaining model on the CPU storage~\cite{aminabadi2022deepspeedinferenceenablingefficient, sheng2023flexgenhighthroughputgenerativeinference}. Researchers have also explored the accuracy vs compute tradeoff space so as to find the right quantization\cite{quant, quant1} method to drop the less useful data and compute while still achieving a satisfactory accuracy. Although these optimizations have been proven useful to cut down the power budget, decrease latency and help in user's data security, these are only useful in generating response from the LLM faster and do not focus on integrating user's context in the generation output. Also, for a problem which presents as many diverse architectural aspects as RAG, the software solution needs to cater to the diverse architectures that could be running them. 

% % \textbf{Hardware resource management:} There have been works in the direction of resource-management for CPUs~\cite{gong2022graphite, \Cref, nair2024parallelization} which attempts to overlap computation and memory accesses, and better use mutltiple cores. However, these works are not directly applicable to RAG application on edge devices. 

% \section{Related Work} 
% \label{sec:related_works} %Cyan
% Recent research and industrial adoption underscore the importance of improving both RAG- and LLM-based applications. Several previous works~\cite{cacheblend, hermes, rago} investigated optimizations for server-scale RAG pipelines, including FlashRAG~\cite{jin2024flashrag}, which offers a modular framework that spans diverse indexing and database choices, PipeRAG~\cite{jiang2024piperag}, TurboRAG~\cite{lu2024turborag}, and others~\cite{zhang2024accelerating,EdgeRAG,ragcacheefficientknowledgecaching} that present system-level insights, often targeting datacenter deployments. Some proposals improve RAG performance through specialized near storage~\cite{RAGAS, reis}, or in-memory computation~\cite{jiang2023chameleon,qin2024robust, HETERag, accelRAG, drex}, delivering notable speedups and energy efficiency, but complicating integration with resource-limited devices. For scenarios where data privacy is paramount, offloading computation to the cloud, as in \cite{huang2024toward}, can introduce confidentiality concerns. Hence, we focus on a platform-flexible software solution that suits edge settings while addressing memory and compute constraints.

% Large LLMs such as LLaMA~\cite{touvron2023llama}, Mixtral~\cite{jiang2024mixtral}, and GPT3~\cite{gpt3} have sparked research into model and resource optimization techniques for edge deployments~\cite{10.1145/emperical}. Various approaches schedule heavy layers on GPUs and offload others to CPU storage~\cite{aminabadi2022deepspeedinferenceenablingefficient, sheng2023flexgenhighthroughputgenerativeinference}, or use carefully tuned quantization schemes~\cite{quant, quant1}, to reduce memory footprint with minimal accuracy loss. Such methods are critical for running large models on low-power devices, but they center on speeding up the core generative model rather than incorporating local context and retrieval. A broader software approach is needed to account for the different stages of the RAG pipeline and the wide variety of device architectures on which they operate.

% Finally, a body of work on hardware resource management for CPUs~\cite{gong2022graphite, nair2024parallelization} explores better concurrency and compute–memory overlap, but these solutions are not directly tailored to the unique combination of encoding, retrieval, and generation of RAG on edge platforms. Our work bridges these gaps by proposing a system-level resource-aware pipeline that addresses challenges ranging from CPU/ GPU partitioning and memory usage to overall QoS in personal computing devices.
